
The discrete solution is lifted into continuous space and refined using clamped cubic B-splines. Let \(P_{\text{disc}}\) denote the discrete path obtained from the joint A* search, and let \(R(\cdot)\) denote the spline refinement operator that maps \(P_{\text{disc}}\) to a continuous trajectory \(\tau_{\text{spline}}\).

We define:
\begin{equation}
\tau_{\mathrm{spline}} = \mathcal{R}(P_{\mathrm{disc}}).
\end{equation}

The refinement operator enforces $C^2$ continuity and curvature feasibility constraints, while producing a smooth trajectory that interpolates or approximates the discrete path. However, this transformation introduces geometric distortion relative to the original discrete solution.

We model this distortion as an additional multiplicative suboptimality term:
\begin{equation}
L(\tau_{\mathrm{spline}}) \le (1 + \epsilon_{\mathrm{spline}})\, L(P_{\mathrm{disc}}),
\end{equation}
where $\epsilon_{\mathrm{spline}} \ge 0$ captures the approximation error induced by spline smoothing, curvature enforcement, and deviation from piecewise-linear geodesics.

Unlike the discrete solution, spline refinement is not cost-preserving in general, as curvature constraints and smoothness requirements may increase trajectory length relative to the discrete path.